% Generated draft: requires mathematical and source review.
\documentclass[11pt]{article}
\usepackage[margin=25mm]{geometry}
\usepackage{fontspec}
\setmainfont{Latin Modern Roman}
\usepackage{amsmath,amssymb,mathtools,unicode-math}
\setmathfont{Latin Modern Math}
\usepackage{enumitem,needspace,xurl,hyperref}
\setlength{\parindent}{0pt}
\setlength{\parskip}{5pt}
\setlength{\emergencystretch}{3em}
\newcommand{\N}{\mathbb{N}}
\newcommand{\Z}{\mathbb{Z}}
\newcommand{\Q}{\mathbb{Q}}
\newcommand{\R}{\mathbb{R}}
\newcommand{\C}{\mathbb{C}}
\newcommand{\F}{\mathbb{F}}
\newcommand{\A}{\mathbb{A}}
\newcommand{\PP}{\mathbb{P}}
\newcommand{\E}{\mathbb{E}}
\newcommand{\eps}{\varepsilon}
\newcommand{\dd}{\,\mathrm{d}}
\begin{document}
\section*{UnsolvedMath ID: 20001658}
\textbf{Source problem code:} catalogue entry
\subsection{Definitions and mathematical statement}
The record merges the end of Section 3, Section 4 and the next heading. Its actual problem is: equip the rapid-decay Sobolev space naturally with operator-space matrix norms, and prove, for example, completely bounded rapid decay for polynomial-growth groups. The neighboring discussion of distorted subgroup lengths is context, not a replacement target.

Let \(\Gamma\) be discrete and let \(L:\Gamma\to[0,\infty)\) satisfy \(L(e)=0\), \(L(g^{-1})=L(g)\), and \(L(gh)\leq L(g)+L(h)\). Let \(\lambda_g\) be the left regular unitaries and \(C_r^*(\Gamma)\) the reduced group \(C^*\)-algebra. The Fourier map on finitely supported functions is
\[
 \mathcal F(f)=\sum_g f(g)\lambda_g.
\]
Ordinary RD with a positive-weight Sobolev convention means bounded extension of this map for some Sobolev exponent. The source's displayed denominator-weighted norm is inconsistent with that inclusion for unbounded \(L\): point masses have norm tending to zero while their Fourier images have norm one. Whether the denominator is read as \(1+L^s\) or \((1+L)^s\) does not cure this sign problem.

Use the aid's explicit positive-weight correction
\[
 H_t=\ell^2(\Gamma,(1+L)^t),\qquad
 \|f\|_{H_t}^2=\sum_g|f(g)|^2(1+L(g))^t.
\]
An operator space has matrix norms satisfying the direct-sum maximum rule and scalar-matrix inequalities
\[
 \|V\oplus W\|=\max\{\|V\|,\|W\|\},\qquad
 \|\alpha V\beta\|\leq\|\alpha\|\,\|V\|\,\|\beta\|.
\]
The multipliers \(\alpha,\beta\) are scalar rectangular complex matrices; the source's extracted coefficient-space notation for them needs this typing correction. A map is completely bounded if all entrywise matrix amplifications have uniformly bounded operator norm. The target \(C_r^*(\Gamma)\) has its natural concrete operator-space norms.

The proposed domain is the row-column intersection
\[
 RC_t(\Gamma,L)=R(H_t)\cap C(H_t).
\]
For finitely supported matrix coefficients \(F_g\in M_m(\C)\), put \(w_g=(1+L(g))^t\). Its matrix norm is explicitly
\[
 \left\|\sum_gF_g\delta_g\right\|_{M_m(RC_t)}
 =\max\left\{
 \left\|\sum_gw_gF_gF_g^*\right\|^{1/2},
 \left\|\sum_gw_gF_g^*F_g\right\|^{1/2}\right\}.
\]
This fixes the row-column structure and exponent normalization, including compatibility with the involution obtained by inversion and adjoint.

The aid reports, for
\[
 Z_L(t)=\sum_{g\in\Gamma}(1+L(g))^{-t}<\infty,
\]
a completely bounded Fourier extension with
\[
 \|\mathcal F:RC_t\to C_r^*(\Gamma)\|_{\mathrm{cb}}
 \leq Z_L(t)^{1/2}.
\]
If \(\Gamma\) is amenable, it reports necessity of the same summability condition and the exact constant
\[
 \|\mathcal F\|_{\mathrm{cb}}^2=Z_L(t).
\]
Thus in the amenable case boundedness already at scalar level, complete boundedness, and convergence of \(Z_L(t)\) are equivalent. Necessity is not asserted here for all nonamenable groups.

In particular, if \(|\{g:L(g)\leq R\}|\leq C(1+R)^D\), every \(t>D\) gives the reported completely bounded RD property. For a finitely generated polynomial-growth group this applies to a word length, with \(D\) a ball-growth exponent. Under the alternative weight \((1+L)^{2s}\), the condition is \(s>D/2\). Polynomial word growth alone does not give the same conclusion for an arbitrary unrelated length.

The neighboring source context says RD passes to subgroups with induced length, mentions polynomial distortion constraints for cyclic subgroups, and gives a free-fiber example \(F_2< F_2\rtimes\Z\) with exponentially distorted induced length. Its extracted presentation \(ac=a2b,bc=ab\) has lost exponent/conjugation typography; a coherent interpretation is an automorphism \(a\mapsto a^2b,b\mapsto ab\), with stable-letter convention needing verification. Those claims are retained as source context, not used to establish the operator-space theorem.
\subsection{Short English statement}
Repair the sign of the extracted Sobolev weight and ask for a natural operator-space form of rapid decay. The aid proposes the row-column Hilbert intersection, gives a summability bound for the completely bounded Fourier map, and reports an exact norm and necessity for amenable groups. Polynomial growth yields the result for the corresponding word-length growth bound.
\subsection{Sources}
\begin{itemize}
\item Local UnsolvedMath registry, permanent ID \(20001658\), code catalogue entry; canonical \texttt{aim-geometric-group-theory-notes.json}, \texttt{notes[166]}.
\item AIM workshop note ``Property of rapid decay,'' 21 March 2006, Section 4 ``Completely bounded property RD,'' Definition 4.1, Remark 4.2 and its final problem paragraph, \url{https://aimath.org/WWN/rapiddecay/rapiddecay.pdf}. The registry's item ``1'' is part of the merged extraction; the supplied verification identifies Section 4 as the actual problem.
\item The supplied record has no additional authoritative paper citation for the row-column construction, sharp amenable norm or distorted-fiber example. No PDF inspection is asserted here.
\item Research attempt 1: ``partial\_result,'' ``partially\_solved,'' rigor label ``proved,'' novelty confidence ``low.'' The machine-generated aid is not peer reviewed; independently verify before citation.
\end{itemize}
\end{document}
